\label{sec:results}


\subsection{Main Results}
\label{sec:main_result}

We use outline coverage as a proxy to assess the pre-writing stage (see~\refsec{sec:task_outline}).~\reftab{table:outline_quality} shows the heading soft recall and entity recall. %
Outlines directly generated by LLMs (\textit{Direct Gen}) already demonstrate high heading soft recall, indicating LLMs' ability to grasp high-level aspects of a topic through their rich parametric knowledge. However, \system, by asking effective questions to research the topic, can create higher recall outlines that cover more topic-specific aspects. Notably, although \textit{RAG} leverages additional information, %
presenting unorganized information in the context window makes outline generation more challenging for the weaker model, \ie, GPT-3.5, leading to worse performance. To test the limit of the \textit{RAG} baseline, we further expand the retrieved sources by starting with the outline produced by \textit{RAG}, using its section titles as search queries to collect more sources, and inputting the newly collected sources together with the initial outline to LLM to generate a polished outline. This modified approach is referred to as ``\textit{RAG-expand}'' in~\reftab{table:outline_quality}. The experiment results indicate that even though having an additional round of search and refinement can improve the outline produced by \textit{RAG}, our proposed \system still surpasses its performance.

We further evaluate the full-length article quality. %
As shown in~\reftab{table:article_quality}, \textit{oRAG} significantly outperforms \textit{RAG}, highlighting the effectiveness of using outlines for structuring full-length article generation. Despite this method's advantages in leveraging retrieval and outlining, our approach still outperforms it. %
The effective question asking mechanism %
enhances the articles with greater entity recall. The evaluator LLM also rates these articles with significantly higher scores in the aspects of ``Interest Level'',  ``Relevance and Focus'', and ``Coverage''. Nonetheless, we acknowledge the possibility of the evaluator LLM overrating machine-generated text. Our careful human evaluation (\refsec{sec:human_eval}) reveals that \system still has much room for improvement.

\begin{table}
\centering
\resizebox{0.8\columnwidth}{!}{%
\begin{tabular}{lcc} 
\toprule
                       & \textbf{Citation Recall} & \textbf{Citation Precision}  \\ 
\midrule
\system & 84.83           & 85.18               \\
\bottomrule
\end{tabular}
}
\caption{Citation quality judged by Mistral 7B-Instruct.}
\label{table:citation_quality}
\vspace{-0.5em}
\end{table}


Although this work primarily focuses on the pre-writing stage and does not optimize generating text with citations, we still examine the citation quality of articles produced by our approach. As reported in~\reftab{table:citation_quality}, Mistral 7B-Instruct judges 84.83\% of the sentences are supported by their citations. Appendix~\ref{appendix:citation_quality_discuss} investigates the unsupported sentences and reveals that the primary issues stem from drawing improper inferences and inaccurate paraphrasing, rather than hallucinating non-existent contents.



\subsection{Ablation Studies}
\begin{table}
\centering
\resizebox{\columnwidth}{!}{%
\begin{tabular}{lccc} 
\toprule
                      & \system & w/o Perspective & w/o Conversation  \\ 
\midrule
$|\mathcal{R}|$ & \textbf{99.83} & 54.36           & 39.56             \\
\bottomrule
\end{tabular}
}
\caption{Average number of unique references ($|\mathcal{R}|$) collected using different methods.}
\label{table:ablation_d}
\end{table}

As introduced in~\refsec{sec:method}, \system prompts LLMs to ask effective questions by discovering specific perspectives and simulating multi-turn conversations. We conduct the ablation study on outline creation by comparing \system with two variants: (1) ``\system w/o Perspective'', which omits perspective in the question generation prompt; (2) ``\system w/o Conversation'', which prompts LLMs to generate a set number of questions altogether. To ensure a fair comparison, we control an equal total number of generated questions across all variants. \reftab{table:outline_quality} shows the ablation results and full \system pipeline produces outlines with the highest recall. Also, ``\system w/o Conversation'' gives much worse results, indicating reading relevant information is crucial to generating effective questions. We further examine how many unique sources are collected in $\mathcal{R}$ via different variants. As shown in~\reftab{table:ablation_d}, the full pipeline discovers more different sources and the trend is in accord with the automatic metrics for outline quality.

We also verify whether having an outline stage is necessary with \system. In~\reftab{table:article_quality}, ``\system w/o Outline Stage'' denotes the results of generating the entire article given the topic and the simulated conversations. %
Removing the outline stage significantly deteriorates the performance across all metrics.










